% Options for packages loaded elsewhere
\PassOptionsToPackage{unicode}{hyperref}
\PassOptionsToPackage{hyphens}{url}
\PassOptionsToPackage{dvipsnames,svgnames,x11names}{xcolor}
%
\documentclass[
  10pt,
]{article}
\usepackage{amsmath,amssymb}
\usepackage{iftex}
\ifPDFTeX
  \usepackage[T1]{fontenc}
  \usepackage[utf8]{inputenc}
  \usepackage{textcomp} % provide euro and other symbols
\else % if luatex or xetex
  \usepackage{unicode-math} % this also loads fontspec
  \defaultfontfeatures{Scale=MatchLowercase}
  \defaultfontfeatures[\rmfamily]{Ligatures=TeX,Scale=1}
\fi
\ifPDFTeX\else
  % xetex/luatex font selection
  \setmainfont[]{DejaVu Serif}
  \setmonofont[]{DejaVu Sans Mono}
\fi
% Use upquote if available, for straight quotes in verbatim environments
\IfFileExists{upquote.sty}{\usepackage{upquote}}{}
\IfFileExists{microtype.sty}{% use microtype if available
  \usepackage[]{microtype}
  \UseMicrotypeSet[protrusion]{basicmath} % disable protrusion for tt fonts
}{}
\makeatletter
\@ifundefined{KOMAClassName}{% if non-KOMA class
  \IfFileExists{parskip.sty}{%
    \usepackage{parskip}
  }{% else
    \setlength{\parindent}{0pt}
    \setlength{\parskip}{6pt plus 2pt minus 1pt}}
}{% if KOMA class
  \KOMAoptions{parskip=half}}
\makeatother
\usepackage{xcolor}
\usepackage[margin=2cm]{geometry}
\usepackage{longtable,booktabs,array}
\usepackage{calc} % for calculating minipage widths
% Correct order of tables after \paragraph or \subparagraph
\usepackage{etoolbox}
\makeatletter
\patchcmd\longtable{\par}{\if@noskipsec\mbox{}\fi\par}{}{}
\makeatother
% Allow footnotes in longtable head/foot
\IfFileExists{footnotehyper.sty}{\usepackage{footnotehyper}}{\usepackage{footnote}}
\makesavenoteenv{longtable}
\setlength{\emergencystretch}{3em} % prevent overfull lines
\providecommand{\tightlist}{%
  \setlength{\itemsep}{0pt}\setlength{\parskip}{0pt}}
\setcounter{secnumdepth}{-\maxdimen} % remove section numbering
\ifLuaTeX
  \usepackage{selnolig}  % disable illegal ligatures
\fi
\IfFileExists{bookmark.sty}{\usepackage{bookmark}}{\usepackage{hyperref}}
\IfFileExists{xurl.sty}{\usepackage{xurl}}{} % add URL line breaks if available
\urlstyle{same}
\hypersetup{
  pdftitle={Code run-book for Claude Code: bc-grade-opendata v1.0},
  colorlinks=true,
  linkcolor={Maroon},
  filecolor={Maroon},
  citecolor={Blue},
  urlcolor={Blue},
  pdfcreator={LaTeX via pandoc}}

\title{Code run-book for Claude Code: bc-grade-opendata v1.0}
\author{}
\date{}

\begin{document}
\maketitle

{
\hypersetup{linkcolor=}
\setcounter{tocdepth}{3}
\tableofcontents
}
\hypertarget{runbook-taking-v0.1-to-a-submission-ready-paper-with-claude-code}{%
\section{RUNBOOK --- taking v0.1 to a submission-ready paper with Claude
Code}\label{runbook-taking-v0.1-to-a-submission-ready-paper-with-claude-code}}

This runbook is written for a Claude Code session (or a human analyst)
working on a branch of this repository. Each step has an
\textbf{action}, the \textbf{files it touches} and an \textbf{acceptance
check}. Work top to bottom; tick the status table as you go and commit
after each step.

\begin{verbatim}
git checkout -b submission-run
claude                      # then: "Read CLAUDE.md and RUNBOOK.md, then do step R1"
\end{verbatim}

\hypertarget{status-at-v1.0-2026-10-18-submission-build}{%
\subsection{Status at v1.0 (2026-10-18, submission
build)}\label{status-at-v1.0-2026-10-18-submission-build}}

\begin{longtable}[]{@{}
  >{\raggedright\arraybackslash}p{(\columnwidth - 2\tabcolsep) * \real{0.5000}}
  >{\raggedright\arraybackslash}p{(\columnwidth - 2\tabcolsep) * \real{0.5000}}@{}}
\toprule\noalign{}
\begin{minipage}[b]{\linewidth}\raggedright
Area
\end{minipage} & \begin{minipage}[b]{\linewidth}\raggedright
State
\end{minipage} \\
\midrule\noalign{}
\endhead
\bottomrule\noalign{}
\endlastfoot
R1 data + guards & Done --- checksums unchanged since v0.1; 12 tests
pass (2 new: MI pooling, Duke in-situ/components) \\
R2 exact reproduction of v0.1 & Done --- \texttt{numbers.tex} identical
after a clean re-download and re-run \\
R3 submission resampling (10 repeats, 2000 bootstraps) & Done \\
R4 multiple imputation (m = 20, pooled predictions) + complete case &
Done --- \texttt{e2mi}, \texttt{tab\_mi} \\
R5 TabPFN & \textbf{Blocked} --- weights gated on Hugging Face (login +
licence) and huggingface.co unreachable from the build sandbox; code
path in \texttt{make\_model("tabpfn")}, runs automatically when weights
are cached (\texttt{optional\_models\_status.json}) \\
R6 more Saudi/Gulf anchors & Partly done --- Riyadh (Durrani 2025),
Najran, Madinah added; pooled random-effects anchor; age-share
sensitivity; bootstrap CIs. \textbf{All new values are
abstract/snippet-level} (journal sites unreachable) → V2 must re-read
the tables \\
R7 sample size (Riley/Pate multinomial) & Done --- \texttt{e10};
criteria \textbf{not met} for the G1-vs-G2 sub-model (reported in the
paper) \\
R8--R9 image-level arm (KAU-BCMD, CDD-CESM, BUSI) & \textbf{Blocked} ---
Kaggle/TCIA hosts unreachable from the build sandbox \\
R10 Duke-Breast-Cancer-MRI & Done --- clinical + 529 radiomics from a
checksummed GitHub mirror of the TCIA XLSX; Duke is a second
full-feature external cohort, in LODO, leakage audit (component scores)
and E9 radiomics fusion \\
R11 manuscript & Done --- rebuilt: Figure 1 framework (TikZ), pgfplots
charts from \texttt{generated/data/*.csv}, new Introduction/Discussion,
abstract \\
V1 citations & Done via search-engine records (79 entries;
\texttt{paper/refs\_annotated.bib} keeps provenance notes). One
DOI-resolver pass from an unrestricted machine still advised \\
V2 anchor double-read & \textbf{To do (human)} \\
V3 TRIPOD+AI full checklist & Done (Supplementary Table) --- STROBE
checklist still to add \\
V4 independent clean re-run & To do on a second machine \\
G1 GitHub + Zenodo & Package ready (\texttt{.github/workflows/ci.yml}
added); repository creation + DOI = PI \\
Author list, affiliations, funding, CRediT, PPI statement & To do PI \\
\end{longtable}

Headline v1.0 results (from \texttt{paper/generated/numbers.tex}):
internal macro-AUC 0.72--0.73 for every learned model, XGBoost − LR =
+0.001 (95\% CI −0.006 to +0.008); external full-feature macro-AUC
0.69--0.71 (MSK-2018), 0.70--0.72 (Duke); grade-3 calibration slope ≈ 1
externally but intercept +0.99 (MSK) and −1.02 (Duke), both fixed by
intercept update; MI vs single imputation Δ ≤ 0.003; leakage inflation
up to +0.30; Duke radiomics fusion −0.045 vs clinical; Riley sample size
not met (4477 required vs 2386); Saudi case-mix raises each cohort's
grade-3 share by ≈ +10 points, and a METABRIC model would predict 59\%
grade 3 vs 36.7\% published (Dammam).

Headline v0.1 results (superseded; see
\texttt{paper/generated/numbers.tex} for v1.0) (all from
\texttt{paper/generated/numbers.tex}): internal macro-AUC ≈ 0.72 for
every learned model (logistic regression = XGBoost); external macro-AUC
0.64--0.80; grade-3 calibration slope ≈ 1.03 in MSK-2018 but intercept ≈
+0.99 (baseline-risk shift, fixed by intercept update); 7/7 SHAP
directions plausible; leakage inflates macro-AUC by up to +0.27; Saudi
re-weighting implies ≈ 61\% grade 3 vs 36.7\% published → Western models
would over-predict grade 3 in Saudi patients.

\begin{center}\rule{0.5\linewidth}{0.5pt}\end{center}

\hypertarget{r0-environment-15-min}{%
\subsection{R0 --- Environment (15 min)}\label{r0-environment-15-min}}

\textbf{Action.}
\texttt{python3\ -m\ venv\ .venv\ \&\&\ .\ .venv/bin/activate\ \&\&\ pip\ install\ -r\ requirements.txt}.
LaTeX: TeX Live with \texttt{charter}, \texttt{natbib},
\texttt{booktabs}, \texttt{enumitem}, \texttt{hyperref},
\texttt{latexmk}. \textbf{Accept.}
\texttt{python\ -c\ "import\ sklearn,\ xgboost,\ lightgbm,\ shap,\ statsmodels,\ imblearn"}
succeeds.

\hypertarget{r1-data-guards-5-min}{%
\subsection{R1 --- Data + guards (5 min)}\label{r1-data-guards-5-min}}

\textbf{Action.} \texttt{make\ data\ test}. \textbf{Accept.} 10 tests
pass. \texttt{git\ diff\ data/raw/MANIFEST.json} is empty; if checksums
changed, cBioPortal updated a study --- record the date in
\texttt{paper/main.tex} Data availability and re-run everything.

\hypertarget{r2-reproduce-v0.1-exactly-15-min-2-cores}{%
\subsection{R2 --- Reproduce v0.1 exactly (≈ 15 min, 2
cores)}\label{r2-reproduce-v0.1-exactly-15-min-2-cores}}

\textbf{Action.} \texttt{make\ experiments\ report\ paper}.
\textbf{Accept.} \texttt{git\ diff\ paper/generated/numbers.tex} shows
no change (the seed is fixed). If numbers move by more than ±0.002,
compare package versions with \texttt{results/run\_*.json} before
continuing.

\hypertarget{r3-submission-grade-resampling}{%
\subsection{R3 --- Submission-grade
resampling}\label{r3-submission-grade-resampling}}

\textbf{Action.} In \texttt{configs/experiment.yaml} set
\texttt{cv.repeats:\ 10}, \texttt{bootstrap:\ 2000}. Run
\texttt{make\ experiments\ report\ paper}. \textbf{Accept.} Methods text
updates automatically through \texttt{\textbackslash{}Repeats} /
\texttt{\textbackslash{}Boot}. Macro-AUCs change \textless{} 0.005.

\hypertarget{r4-multiple-imputation-with-pooled-predictions}{%
\subsection{R4 --- Multiple imputation with pooled
predictions}\label{r4-multiple-imputation-with-pooled-predictions}}

\textbf{Why.} Reviewers will ask; v0.1 uses single in-fold
chained-equation imputation. \textbf{Action.} In
\texttt{src/bcgrade/models.py} add \texttt{MIClassifier(m=20)}: inside
\texttt{fit}, create m
\texttt{IterativeImputer(sample\_posterior=True,\ random\_state=seed+i)}
→ scaler → clone(model) pipelines; in \texttt{predict\_proba}, average
the m probability matrices. Register as model option \texttt{mi=True};
add a test that predictions are averaged over m distinct imputations.
Add an experiment \texttt{e2mi} for \texttt{mlr} and \texttt{xgb} only.
Add a complete-case sensitivity row. Report in a new supplementary table
\texttt{tab\_mi.tex}. \textbf{Accept.} MI and single-imputation
macro-AUC differ by \textless{} 0.01, or the paper discusses why not.

\hypertarget{r5-tabular-foundation-model-tabpfn}{%
\subsection{R5 --- Tabular foundation model
(TabPFN)}\label{r5-tabular-foundation-model-tabpfn}}

\textbf{Action.} \texttt{pip\ install\ tabpfn} (needs model weights from
Hugging Face on first use). Add \texttt{"tabpfn"} to
\texttt{make\_model} (no tuning grid;
\texttt{TabPFNClassifier(random\_state=seed)}), include in E2/E3/LODO.
Cite \texttt{hollmann2025tabpfn} (already in refs.bib). \textbf{Accept.}
Runs within the CV harness without leakage (it sits after impute/scale
in the Pipeline). Expectation from v0.1: it will not beat regression by
more than the CI; report either way.

\hypertarget{r6-strengthen-the-saudi-gulf-anchor-key-for-sponsors}{%
\subsection{R6 --- Strengthen the Saudi / Gulf anchor (key for
sponsors)}\label{r6-strengthen-the-saudi-gulf-anchor-key-for-sponsors}}

\textbf{Action.} 1. Find further published Saudi/Gulf
clinicopathological series reporting grade, ER/PR/HER2, histology and
age (Saudi Medical Journal, Annals of Saudi Medicine, J Clin Med,
Cureus, BMC Cancer; KFSH\&RC Riyadh, KAMC Riyadh/Jeddah, KAUH Jeddah;
also UAE, Kuwait, Oman, Qatar). Transcribe Table 1 counts into
\texttt{configs/anchors.yaml} (one block per study, with page and table
numbers). 2. Generalise \texttt{e7} to loop over all anchors and add a
random-effects pooled Saudi grade distribution. 3. Add a ``joint''
sensitivity: rake on ER × HER2 jointly if any study reports the
cross-tab. 4. Add Saudi Cancer Registry figures used in the Introduction
(Basudan 2022 already; add the most recent SCR report). \textbf{Accept.}
Every anchor value has a citation + table/page reference in a YAML
comment and is double-checked by a second reader (V2).

\hypertarget{r7-sample-size-justification}{%
\subsection{R7 --- Sample-size
justification}\label{r7-sample-size-justification}}

\textbf{Action.} Use \texttt{pmsampsize} (R) or a Python port for a
multinomial outcome (Riley et al.~BMJ 2020; Pate et al. for
multinomial). Inputs: 10 parameters, smallest class prevalence 0.089,
anticipated Cox--Snell R² from v0.1 (compute from METABRIC logistic
model). Write the output into \texttt{results/sample\_size.json} and a
sentence in Methods via a macro. \textbf{Accept.} Required n ≤ 2386, or
the paper says it is not met.

\hypertarget{r8-imaging-arm-b1-kau-bcmd-saudi-mammography}{%
\subsection{R8 --- Imaging arm B1: KAU-BCMD (Saudi
mammography)}\label{r8-imaging-arm-b1-kau-bcmd-saudi-mammography}}

\textbf{Action.} 1. \texttt{pip\ install\ kaggle}; place
\texttt{\textasciitilde{}/.kaggle/kaggle.json};
\texttt{python\ scripts/download\_imaging.py\ kau\_bcmd}. 2. Build
\texttt{src/bcgrade/imaging/kau.py}: index images → patient/case id →
BI-RADS label; record the version (the Kaggle folders hold BI-RADS 1, 3,
4, 5; the paper describes 1--5 --- document what you downloaded). 3.
Task: BI-RADS 1 vs 3 vs 4--5 (or benign-appearing vs suspicious),
\textbf{StratifiedGroupKFold by case}. 4. Models: frozen ImageNet
ViT/ConvNeXt or a mammography foundation model → embeddings → logistic
regression; then fine-tuned EfficientNet-B0 on a GPU. Same
metrics/calibration code. 5. Leakage demo: image-level vs case-level
split (report like E4-d). \textbf{Accept.} No case appears in two folds
(unit test). Results in a new section ``Imaging arm''.

\hypertarget{r9-imaging-arm-b2-cdd-cesm-egypt-and-busi-egypt}{%
\subsection{R9 --- Imaging arm B2: CDD-CESM (Egypt) and BUSI
(Egypt)}\label{r9-imaging-arm-b2-cdd-cesm-egypt-and-busi-egypt}}

\textbf{Action.}
\texttt{python\ scripts/download\_imaging.py\ cdd\_cesm\ -\/-manifest-only},
then full download; BUSI via Kaggle. Tasks: pathology benign vs
malignant (CDD-CESM, BUSI) and BI-RADS (CDD-CESM). Cross-site check:
train on CDD-CESM low-energy images, test on KAU-BCMD for a ``suspicious
vs not'' label mapping (document the mapping table in
\texttt{configs/}). \textbf{Accept.} Patient-grouped CV; external
cross-site AUC reported with CI.

\hypertarget{r10-duke-breast-cancer-mri-clinical-radiomics-restores-bi-rads-era-variables}{%
\subsection{R10 --- Duke-Breast-Cancer-MRI clinical + radiomics
(restores BI-RADS-era
variables)}\label{r10-duke-breast-cancer-mri-clinical-radiomics-restores-bi-rads-era-variables}}

\textbf{Why.} It is the only open cohort linking Nottingham grade to
mammographic calcification, density and MRI radiomics (922 patients; CC
BY-NC 4.0). \textbf{Action.} Download the clinical XLSX and the
529-feature radiomics XLSX from TCIA (no image download needed). Add
\texttt{duke()} to \texttt{harmonise.py} (map Nottingham grade; drop
tubule/nuclear/mitotic scores --- leakage; code calcification and
density). Add Duke as (a) external validation cohort for the full CDM
model, (b) a clinical + radiomics fusion model (E9), compared with the
published radiomics benchmark (validation AUC 0.722 for grade 3 vs 1--2;
cite after V1 verification). \textbf{Accept.} New external row for Duke;
licence note (non-commercial) added to Data availability.

\hypertarget{r11-manuscript-update}{%
\subsection{R11 --- Manuscript update}\label{r11-manuscript-update}}

\textbf{Action.} Re-run \texttt{make\ report\ paper}. Rewrite only
prose; never hard-code a number (CLAUDE.md rule 1). Update the abstract
(≤ 300 words for most journals) and the TRIPOD+AI table.
\textbf{Accept.}
\texttt{grep\ -nE\ "{[}0-9{]}\textbackslash{}.{[}0-9{]}\{2,\}"\ paper/main.tex\ paper/sections/*.tex}
shows only published anchor figures (Saudi/registry) and software
versions, no results.

\hypertarget{v1-citation-verification}{%
\subsection{V1 --- Citation
verification}\label{v1-citation-verification}}

Check every \texttt{refs.bib} entry with a DOI resolver; fix all
\texttt{VERIFY} notes (wang2022duke, kim2025grade, kussaibi2021saudi
DOI, mbcproject2022). Remove uncited entries.

\hypertarget{v2-anchor-verification}{%
\subsection{V2 --- Anchor verification}\label{v2-anchor-verification}}

Second reader re-transcribes \texttt{configs/anchors.yaml} from the
source PDFs blind, then diff.

\hypertarget{v3-reporting-checklists}{%
\subsection{V3 --- Reporting checklists}\label{v3-reporting-checklists}}

Complete the full TRIPOD+AI (27 items) and STROBE checklists as
supplementary PDFs; replace the abbreviated Table S5.

\hypertarget{v4-independent-re-run}{%
\subsection{V4 --- Independent re-run}\label{v4-independent-re-run}}

On a clean machine: \texttt{git\ clone\ …\ \&\&\ make\ all};
\texttt{numbers.tex} must match the committed version.

\hypertarget{g1-github-zenodo-release}{%
\subsection{G1 --- GitHub + Zenodo
release}\label{g1-github-zenodo-release}}

\begin{enumerate}
\def\labelenumi{\arabic{enumi}.}
\tightlist
\item
  Create the GitHub repository (public), push \texttt{main} (data and
  \texttt{results/*.npz} are git-ignored).
\item
  Enable the Zenodo--GitHub integration, tag \texttt{v1.0.0}, cite the
  DOI in \texttt{paper/main.tex} and \texttt{CITATION.cff}.
\item
  Add a GitHub Action running \texttt{make\ data\ test} on push (CPU,
  \textless{} 5 min).
\end{enumerate}

\begin{center}\rule{0.5\linewidth}{0.5pt}\end{center}

\hypertarget{journal-targeting-note}{%
\subsubsection{Journal targeting note}\label{journal-targeting-note}}

Q1 options that publish open-data prediction studies with TRIPOD+AI:
\emph{Breast Cancer Research and Treatment}, \emph{The Breast},
\emph{JCO Clinical Cancer Informatics}, \emph{npj Breast Cancer} (needs
imaging/genomic angle --- R8--R10), \emph{Diagnostic and Prognostic
Research} (methods angle: the leakage audit). Decide after R10.

\hypertarget{claude.md-working-rules-for-claude-code-in-this-repository}{%
\section{CLAUDE.md --- working rules for Claude Code in this
repository}\label{claude.md-working-rules-for-claude-code-in-this-repository}}

This repo is the reproducibility package for the manuscript in
\texttt{paper/main.tex} (open-data breast cancer histological-grade
prediction). Start every session by reading \texttt{RUNBOOK.md} and the
``Status'' table at its top.

\hypertarget{non-negotiable-rules}{%
\subsection{Non-negotiable rules}\label{non-negotiable-rules}}

\begin{enumerate}
\def\labelenumi{\arabic{enumi}.}
\tightlist
\item
  \textbf{No number is typed by hand into the paper.} Every figure,
  table and in-text number comes from \texttt{results/} →
  \texttt{experiments/report.py} → \texttt{paper/generated/}
  (\texttt{numbers.tex} macros, \texttt{tab\_*.tex},
  \texttt{fig\_*.pdf}). If you change an analysis, re-run
  \texttt{make\ experiments\ report\ paper} and read the diff of
  \texttt{numbers.tex}.
\item
  \textbf{Every learned step stays inside the sklearn/imblearn Pipeline}
  (array copy → imputation → scaling → optional SMOTE → model). Never
  fit an imputer, scaler, resampler or feature selector on data that
  includes a test fold. \texttt{tests/test\_pipeline.py} guards this;
  add a test when you add a step.
\item
  \textbf{Leakage columns are never predictors} of grade: NPI (contains
  grade), Nottingham component scores (tubule/pleomorphism/mitotic),
  stage (contains T/N but also used for selection),
  PAM50/IntClust/Oncotype unless explicitly in a sensitivity analysis.
  Keep \texttt{test\_leakage\_columns\_never\_in\_feature\_sets}
  passing.
\item
  \textbf{In-situ disease is excluded, never relabelled.} ILC is never
  merged with IDC. Grade is invasive Nottingham/differentiation grade
  1--3 only.
\item
  \textbf{One row per patient.} Any image or multi-sample source must be
  split with patient-grouped CV (\texttt{StratifiedGroupKFold}).
\item
  \textbf{Data are not redistributed.} \texttt{data/} is git-ignored.
  Download scripts + \texttt{data/raw/MANIFEST.json} (SHA-256) are the
  provenance record. Respect CC BY-NC terms (Duke MRI, BRACS). Never use
  a source that needs a signed DUA (VinDr-Mammo, full EMBED) without the
  PI's written decision.
\item
  \textbf{Seeds}: \texttt{configs/experiment.yaml:\ seed}. Do not change
  the seed to improve a result.
\item
  Commit messages: imperative, one logical change each; never commit
  \texttt{data/} or large \texttt{results/*.npz}.
\end{enumerate}

\hypertarget{layout}{%
\subsection{Layout}\label{layout}}

\begin{verbatim}
configs/      datasets.yaml (sources, licences, roles) · experiment.yaml (CV, seeds, raking targets) · anchors.yaml (Saudi published figures)
scripts/      download_tabular.py (open cBioPortal + UCI mirror) · download_imaging.py (arm B, Kaggle/TCIA)
src/bcgrade/  harmonise.py (common data model) · describe.py (scale-appropriate stats) · models.py (ladder) · evaluate.py (metrics, CI, calibration, DCA, raking)
experiments/  run.py (E1–E8) · report.py (tables/figures/macros)
tests/        pytest guards for harmonisation and leakage
paper/        main.tex · refs.bib · generated/ (never edit by hand)
\end{verbatim}

\hypertarget{common-commands}{%
\subsection{Common commands}\label{common-commands}}

\begin{verbatim}
make data          # download + checksum
make test          # guards
make experiments   # E1–E8 (~15–25 min on 2 CPU cores at repeats=5)
make report paper  # regenerate tables/figures/macros and compile the PDF
python experiments/run.py e3 e5   # re-run selected stages (E5 and E7 depend on E2 + E3 outputs)
\end{verbatim}

\end{document}
